% SHARD-ID: AQM-69-QSA-FINITE-SET-DIRECT-SUM-PARTICLE
% SHARD-TITLE: Structureless Finite-Set Direct-Sum Particle Association
% SHARD-SUMMARY: Defines the finite-set direct-sum workflow after adding one internal finite-dimensional complex carrier or Hilbert space per site.
% SHARD-SUMMARY: Forms the one-particle-at-one-of-many-sites carrier as a finite direct sum and works the empty and one-point cases.
% SHARD-SUMMARY: Distinguishes this workflow from tensor-site coexistence and compares the one-dimensional internal-space case with AQM-67 only under stated choices.
% SHARD-KEYWORDS: quantum-system association, finite set, direct sum, one particle, internal carrier, Hilbert caveat

\section{Structureless Finite-Set Direct-Sum Particle Association}
\label{sec:qsa-finite-set-direct-sum-particle}

\subsection{Status and Local Anchors}

This is QSA Step 05 from
\path{docs/quantum_system_associations/PLAN.md}.  It is a local
finite-dimensional construction using \texttt{CONVENTIONS.md} conventions
\texttt{(ap)}, \texttt{(aq)}, \texttt{(ar)}, and \texttt{(as)}, with
AQM-64--AQM-66 as the planning and catalogue anchors.  AQM-67 is the formal
finite-set carrier comparison point, and AQM-68 is the tensor-site comparison
point.

The input is not a bare finite set alone.  A bare finite set \(X\) supplies
only labels for alternatives.  The direct-sum particle workflow starts after
one adds an internal finite-dimensional complex carrier, or a finite-dimensional
complex Hilbert space, at each site.  This shard uses no canonical boson or
fermion field claim from AQM-19.  AQM-23 is used only as a cautionary local
anchor: direct sums in label bookkeeping and Hilbert-space direct sums must
not be interchanged with tensor-product residue-qudit systems without an
explicit workflow.

\subsection{Input Data}

There are two finite-dimensional input versions.

\begin{description}
\item[Carrier-site input.]
For every \(x\in X\), choose a finite-dimensional complex vector space
\(K_x\).  The data are the labelled family \(\{K_x\}_{x\in X}\), not merely
the set \(X\).

\item[Hilbert-site input.]
For every \(x\in X\), choose a finite-dimensional complex Hilbert space
\(\mathcal H_x\).  This includes the inner product on each internal carrier.
\end{description}

The word ``internal'' only records what is carried while the alternative site
label is \(x\).  It does not add dynamics, hopping amplitudes, a Hamiltonian,
a preferred observable algebra, or a tensor-factor interpretation.

\subsection{Direct-Sum Carrier}

Under convention \texttt{(as)}, the carrier-site version is
\[
  K_{\oplus}(X)=\bigoplus_{x\in X}K_x .
\]
An element is a tuple \(\psi=(\psi_x)_{x\in X}\) with
\(\psi_x\in K_x\).  Since \(X\) is finite, no support or completion issue is
present.  The canonical injections and projections are denoted
\[
  j_x:K_x\longrightarrow K_{\oplus}(X),
  \qquad
  p_x:K_{\oplus}(X)\longrightarrow K_x,
\]
with \(p_xj_y=0\) for \(x\neq y\) and \(p_xj_x=\operatorname{id}_{K_x}\).
If a displayed listing \(X=\{x_1,\ldots,x_n\}\) is chosen and
\(\dim_{\mathbb C}K_{x_i}=d_i\), then
\[
  \dim_{\mathbb C}K_{\oplus}(X)=\sum_{i=1}^n d_i .
\]
This dimension count is finite direct-sum bookkeeping for the chosen carriers,
not a number determined by the bare set.

In the Hilbert-site version,
\[
  \mathcal H_{\oplus}(X)=\bigoplus_{x\in X}\mathcal H_x
\]
with inner product
\[
  \langle \psi,\phi\rangle_{\oplus}
  =
  \sum_{x\in X}\langle \psi_x,\phi_x\rangle_{\mathcal H_x}.
\]
The orthogonality of distinct summands belongs to this chosen Hilbert direct
sum.  It is not inherited from the bare finite set and was not supplied by
AQM-67.

\subsection{Edge Cases}

\paragraph{Empty set.}
For \(X=\emptyset\),
\[
  K_{\oplus}(\emptyset)=0,
  \qquad
  \mathcal H_{\oplus}(\emptyset)=0
\]
in the corresponding versions.  There are no site labels and no internal
carriers to add.  This differs from the empty tensor-site convention in
AQM-68, where the empty tensor product is \(\mathbb C\).

\paragraph{One-point set.}
For \(X=\{\ast\}\),
\[
  K_{\oplus}(X)=K_\ast,
  \qquad
  \mathcal H_{\oplus}(X)=\mathcal H_\ast .
\]
Thus the one-point direct-sum workflow returns exactly the chosen internal
carrier at the single site.  A one-dimensional result appears only if that
internal carrier was chosen to be one-dimensional.

\subsection{Boundary With Tensor-Site Coexistence}

AQM-68 records the tensor-site workflow
\[
  (X,\{\mathcal H_x\}_{x\in X})
  \longmapsto
  \bigotimes_{x\in X}\mathcal H_x .
\]
That construction represents independent coexisting site factors.  The present
workflow represents alternatives by
\[
  (X,\{\mathcal H_x\}_{x\in X})
  \longmapsto
  \bigoplus_{x\in X}\mathcal H_x .
\]
For displayed finite dimensions \(d_i=\dim\mathcal H_{x_i}\), the tensor-site
carrier has dimension \(\prod_i d_i\), while the direct-sum carrier has
dimension \(\sum_i d_i\).  The empty cases also differ:
\[
  \bigotimes_{x\in\emptyset}\mathcal H_x=\mathbb C,
  \qquad
  \bigoplus_{x\in\emptyset}\mathcal H_x=0.
\]
These are different workflows even when the same labelled family of site
Hilbert spaces is available.

This shard also does not select an observable algebra.  If a later step needs
one, the full algebra
\(\operatorname{End}_{\mathbb C}(\bigoplus_x\mathcal H_x)\) and the
block-preserving algebra \(\bigoplus_x\operatorname{End}_{\mathbb C}(\mathcal
H_x)\) are different choices and must be named.  The tensor-site algebra from
AQM-68 is instead \(\bigotimes_x\operatorname{End}_{\mathbb C}(\mathcal H_x)\).

\subsection{One-Dimensional Internal Spaces and AQM-67}

The comparison with AQM-67 is available only after extra choices are stated.
Assume that every carrier \(K_x\) is a one-dimensional complex vector space
and choose a nonzero vector \(u_x\in K_x\) for each \(x\in X\).  Then
\[
  \Phi:\mathbb C\langle X\rangle
  \longrightarrow
  \bigoplus_{x\in X}K_x,
  \qquad
  \Phi(e_x)=j_x(u_x),
\]
is a vector-space isomorphism.  In the empty case this is the unique map
\(0\to0\).  In the one-point case it is the chosen identification
\(\mathbb C e_\ast\cong K_\ast\).

Without the chosen vectors \(u_x\), the bare set does not identify the formal
generator \(e_x\) with a vector in \(K_x\).  If the \(K_x\) are Hilbert lines,
then chosen unit vectors \(u_x\) make \(\Phi\) compatible with the identity
Gram convention on \(\mathbb C\langle X\rangle\).  That Hilbert comparison is
still a comparison of enhanced data: AQM-67 did not supply inner products,
unit vectors, or orthonormality from the set alone.

\subsection{Conclusion}

The QSA Step 05 output is the alternatives carrier
\[
  \bigoplus_{x\in X}K_x
  \quad\text{or}\quad
  \bigoplus_{x\in X}\mathcal H_x
\]
after the internal carriers or Hilbert spaces have been supplied.  It is not
the formal bare-set carrier of AQM-67 unless one-dimensional internal carriers
and basis vectors are chosen, and it is not the tensor-site coexistence system
of AQM-68.
